\documentclass[12pt]{article}
\usepackage[T1]{fontenc}
\usepackage[utf8]{inputenc}

% Master shell for: (1) Tutorial, (2) Poster, (3) Math Primer
% - One-column, 12pt
% - Uses \input{} so each component remains editable as a separate file.

\usepackage[margin=1in]{geometry}
\usepackage{amsmath, amssymb, mathtools}
\usepackage{enumitem}
\usepackage[hidelinks]{hyperref}
\usepackage{lmodern}

\setlist[itemize]{leftmargin=1.5em}
\setlength{\parskip}{0.6em}
\setlength{\parindent}{0pt}

\title{Appendix A Companion Pack (Tutorial + Poster + Math Primer)}
\date{}

\begin{document}
\maketitle

\section*{Introduction (How to Use This Pack)}
This master file compiles three companion components for Appendix A of the FRFP paper.

\begin{itemize}
\item \textbf{Tutorial (step-by-step)}: a slow walkthrough that tracks clarifications (“tutorial notes”) and ties each concept back to Appendix A.
\item \textbf{Poster (structure-first)}: a block-based, detail-preserving poster-style reorganization for fast review and recall.
\item \textbf{Math Primer (minimum background)}: a poster-style primer covering only the math needed to read Appendix A (objects, morphisms, categories, groupoids, fibers, fibrations, cartesian lifting, etc.).
\end{itemize}

%\paragraph{Compilation note.}
%Place these files in the same directory:
%\begin{itemize}
%\item \texttt{tutorial\_main.tex} (you will supply this)
%\item \texttt{appendixA\_poster\_include.tex} (generated)
%\item \texttt{appendixA\_math\_primer\_include.tex} (generated)
%\end{itemize}

\newpage
\tableofcontents
\newpage

\part*{Tutorial Walkthrough}
\addcontentsline{toc}{part}{Tutorial Walkthrough}
% This file is intended to be included via \input{}.
% Preamble and document environment stripped.

\section*{Tutorial Walkthrough}

\newpage

\section*{How to Use This Companion}

This document is meant to be read \emph{side-by-side} with your original
\textbf{Appendix A --- Mathematical Foundations of FRFP}

\textbf{Important:} 
This companion:

\begin{itemize}
 \item \textbf{does not restate} the formal definitions, theorems, or proofs from Appendix~A;
 \item instead, it \textbf{references} them by their numbering (e.g., ``Definition~1.1''),
 \item and then provides explanations, examples, and prompts that map onto them.
\end{itemize}

In other words, this is a \emph{commentary layer}:

\begin{itemize}
 \item Your original Appendix A remains the source of truth for the math.
 \item This companion explains, illustrates, and contextualizes that math.
\end{itemize}

\subsection*{Pronunciation and Abbreviations}

To make this usable:

\begin{itemize}
 \item The first time a symbol appears, we give its pronunciation in parentheses, e.g.:
 \begin{itemize}
 \item $C_{\text{full}}$ (pronounced \emph{C-full}),
 \item $RB$ (pronounced \emph{R-B}),
 \item $\omega_0$ (pronounced \emph{omega-naught}).
 \end{itemize}
 \item The first time an abbreviation appears, we expand it, e.g.:
 \begin{itemize}
 \item FRFP (Foundational Reasoning and Feedback Protocol),
 \item TDG (Task Decomposition Grammar),
 \item DHO (Damped Harmonic Oscillator).
 \end{itemize}
\end{itemize}

\subsection*{The Running Example: Damped Harmonic Oscillator}

Throughout, we use the damped harmonic oscillator (DHO, pronounced \emph{D-H-O}) as a concrete
example:
\[
\ddot{x}(t) + 2\gamma \dot{x}(t) + \omega_0^2 x(t) = 0,
\]
where:
\begin{itemize}
 \item $x(t)$ is \emph{x-of-t}, the displacement,
 \item $\gamma$ (gamma) is the damping coefficient,
 \item $\omega_0$ (omega-naught) is the natural frequency.
\end{itemize}

For each major definition or construction in Appendix~A, we ask:

\begin{enumerate}
 \item What does this mean in plain language?
 \item How does it apply to the DHO?
 \item How would a system built on top of a ChatGPT-like model prompt or be prompted here?
\end{enumerate}

\newpage

\section{A.0 Foundations (Companion)}

This section comments on \textbf{Section A.0} of Appendix~A, which includes:

\begin{itemize}
 \item A.0.1 Ambient Category $C_{\text{full}}$ (C-full),
 \item A.0.2 Explicit--Tacit Separation (ETS),
 \item A.0.2 (Modal Interpretation of Explicit and Tacit Knowledge),
 \item A.0.3 Explicit Category $E$ (E),
 \item A.0.4 Tacit Category $T$ (T),
 \item A.0.5 Boundary Morphism and $RB$ Functor,
 \item A.0.6 Phase-0 and Phase-1 Axioms.
\end{itemize}

In each subsection below, ``see Definition~X.Y'' refers to your original Appendix A.

\subsection{A.0.1 Ambient Category $C_{\text{full}}$ (Companion)}

\paragraph{Reference in Appendix A.}
See Definition~1.1 (Ambient category).

\subsubsection*{Plain-language view}

You introduce an \emph{ambient category} $C_{\text{full}}$ (pronounced \emph{C-full}) which is the 
mathematical universe where everything lives. Its key objects are:

\begin{itemize}
 \item $\emptyset$ (empty-set): the empty explicit state (no content yet).
 \item $E$ (E): the explicit state object (machine-manipulable representations).
 \item $T$ (T): the tacit state object (human correctness-bearing states).
\end{itemize}

There are morphisms (arrows) between objects, and you group them into three kinds:

\begin{itemize}
 \item \textbf{Explicit morphisms} with codomain $E$: things the AI system does to explicit data.
 \item \textbf{Tacit morphisms} with codomain $T$: things that happen in human cognition.
 \item \textbf{Boundary morphisms} $E \to T$: ways explicit content reaches the human.
\end{itemize}

Among the boundary morphisms, you single out:
\[
RB : E \to T
\]
(Representational Backflow, pronounced \emph{R-B}), modeling the act of presenting representations to
the user.

The category axioms (associativity and identity) are standard: composing transformations is 
well-behaved.

\subsubsection*{DHO Example: what are $E$ and $T$ concretely?}

For the damped harmonic oscillator:
\[
\ddot{x}(t) + 2\gamma \dot{x}(t) + \omega_0^2 x(t) = 0,
\]
explicit states in $E$ might be:

\begin{itemize}
 \item The equation itself, formatted as text or LaTeX.
 \item A state-space representation:
 \[
 \dot{z} =
 \begin{pmatrix}
 0 & 1 \\
 -\omega_0^2 & -2\gamma
 \end{pmatrix} z.
 \]
 \item Numerical solution data: a time grid and values $x(t_i)$.
 \item A symbolic solution formula.
\end{itemize}

Tacit states in $T$ capture things in your head:

\begin{itemize}
 \item ``This looks underdamped.''
 \item ``The decay is too slow; that damping value seems off.''
 \item ``I trust this derivation / I do not trust this derivation.''
\end{itemize}

The boundary morphism $RB : E \to T$ is what happens when the AI displays such representations to
you and you form a tacit opinion about them.

\subsubsection*{Example prompts implementing this structure}

\paragraph{Explicit morphism ($E \to E$).}
\begin{quote}
\textbf{Prompt:} 
``Rewrite the damped harmonic oscillator
$\ddot{x}(t) + 2\gamma \dot{x}(t) + \omega_0^2 x(t) = 0$
as a first-order state-space system. Perform only explicit algebraic rewriting, do not comment on correctness or plausibility.''
\end{quote}

\paragraph{Boundary morphism ($RB : E \to T$).}
The system computes a solution and shows a plot:

\begin{quote}
\textbf{System (explicit output):} 
``Here is the solution for $\gamma = 0.2$, $\omega_0 = 1.0$ and its time-domain plot.'' 
\end{quote}

What happens next (your interpretation) is $RB$ feeding that explicit state into $T$.

\paragraph{Human generating new explicit content (leads to $RI$).}
\begin{quote}
\textbf{User:} 
``That looks underdamped. Please compute the critical damping case $\gamma = \omega_0$ and show the solution.'' 
\end{quote}

This new instruction is a new explicit representation entering $E$ via $RI : \emptyset \to E$.

\subsubsection*{Tracking notes}

\begin{itemize}
 \item $C_{\text{full}}$ (C-full) is the overarching category.
 \item Objects: $\emptyset$, $E$, $T$.
 \item Morphism classes: explicit, tacit, boundary.
 \item Distinguished boundary morphism: $RB : E \to T$.
\end{itemize}

\newpage

\subsection{A.0.2 Explicit--Tacit Separation (ETS) (Companion)}

\paragraph{Reference in Appendix A.}
See Definition~1.2 (Explicit--Tacit Separation (ETS)).

\subsubsection*{Plain-language view}

ETS (Explicit--Tacit Separation, pronounced \emph{E-T-S}) is the axiom that enforces a hard boundary
between AI reasoning and human judgment. It says:

\begin{itemize}
 \item There are \emph{no morphisms} $T \to E$ in $C_{\text{full}}$:
 the AI cannot read your mind or turn tacit states into explicit content on its own.
 \item Tacit evaluation cannot be hidden inside explicit pipelines:
 you cannot smuggle in ``human-like correctness or intuition'' into $EC$ or $ED$.
 \item The boundary morphism $RB : E \to T$ is not invertible:
 once explicit content has been seen and processed by the human, the AI cannot reconstruct
 the tacit state from the explicit trace alone.
\end{itemize}

Conceptually: the AI deals only with representations; humans deal with correctness.

\subsubsection*{DHO Example}

For the DHO, suppose the AI solves:
\[
\ddot{x}(t) + 2\gamma \dot{x}(t) + \omega_0^2 x(t) = 0
\]
and produces the solution
\[
x(t) = e^{-\gamma t}(C_1 \cos(\omega t) + C_2 \sin(\omega t)).
\]

ETS implies:

\begin{itemize}
 \item The AI can check by differentiation and substitution that this satisfies the ODE.
 \item The AI cannot decide whether a human physicist would ``approve'' this solution.
 \item The AI cannot internally simulate a tacit judgement like ``this is physically okay''.
\end{itemize}

Also, after you see the solution, your reaction (trust, doubt, confusion) is in $T$. 
The AI has no morphism $T \to E$ to directly mine that; it only gets new explicit information when
you write something.

\subsubsection*{Prompts that respect / violate ETS}

\paragraph{Respecting ETS.}
\begin{quote}
\textbf{Prompt:} 
``Verify explicitly that the proposed solution satisfies the equation
$\ddot{x} + 2\gamma \dot{x} + \omega_0^2 x = 0$ by differentiation and substitution. 
Do not comment on whether the model is physically appropriate.''
\end{quote}

\paragraph{Violating ETS (what the system must refuse to do).}
\begin{quote}
\textbf{Prompt:} 
``Tell me whether a typical expert would find this solution convincing.'' 
\end{quote}

An FRFP-constrained system should respond:

\begin{quote}
\textbf{System:} 
``I cannot perform tacit evaluation or predict a human expert's correctness judgment. 
I can only check explicit mathematical properties. 
Please specify explicit criteria you want me to verify.''
\end{quote}

\subsubsection*{Tracking notes}

\begin{itemize}
 \item No morphisms $T \to E$.
 \item No tacit evaluation inside explicit pipelines.
 \item $RB$ is not invertible.
\end{itemize}

\newpage

\subsection{A.0.2 Modal Interpretation of Explicit and Tacit Knowledge (Companion)}

\paragraph{Reference in Appendix A.}
See the ``Modal Interpretation'' subsection (Definition~1.3, Definition~1.4 in your text).

\subsubsection*{Plain-language view}

Here you give a modal-logic reading of $E$ and $T$:

\begin{itemize}
 \item $E$ corresponds to what is explicitly represented: formulas, derivations, data.
 \item $T$ corresponds to what is tacitly held to be correct, trusted, accepted.
\end{itemize}

Roughly, explicit content is like ``things written on the board'' and tacit content is like
``things the human believes or endorses''.

Morphisms into $T$ are about updating or refining this correctness-laden state:
``I now think this is correct'', 
``I now reject this argument'', 
and so on.

\subsubsection*{DHO Example}

For the DHO, you might have:

\begin{itemize}
 \item Explicit content in $E$: 
 ``Under the assumption $\gamma < \omega_0$, the system is underdamped and the solution oscillates.''
 \item Tacit content in $T$: 
 ``I (the human) believe this underdamped classification is correct for the particular parameters.''
\end{itemize}

You can have explicit content that you do \emph{not} yet tacitly endorse (e.g., an AI-proposed
solution you are still checking), and you can have tacit beliefs that have not yet been articulated
explicitly.

\subsubsection*{FRFP-style prompting}

A system respecting this modal view might ask you explicitly:

\begin{quote}
\textbf{System:} 
``Here is the derived classification of the damping regime. 
Do you endorse this classification as correct for your physical scenario? 
If not, please specify which part you reject or want to adjust.''
\end{quote}

This invites you to update your tacit state $T$ and, optionally, to produce new explicit content via
$RI$.

\subsubsection*{Tracking notes}

\begin{itemize}
 \item $E$ vs.\ $T$ can be seen as explicit vs.\ correctness-bearing (modal) knowledge.
 \item Tacit updates correspond to morphisms inside $T$ after $RB$ is applied.
\end{itemize}

\newpage

\subsection{A.0.3 Explicit Category $E$ (Companion)}

\paragraph{Reference in Appendix A.}
See Definition~1.5 (Explicit category) and Remark~1.6.

\subsubsection*{Plain-language view}

You define the explicit category $E$ as the part of $C_{\text{full}}$ whose morphisms land in $E$.
Its generators:

\begin{itemize}
 \item $RI : \emptyset \to E$ (Representation Initiation, pronounced \emph{R-I}),
 \item $EC : E \to E$ (Explicit Computation, \emph{E-C}),
 \item $ED : E \to E$ (Explicit Diagnostics, \emph{E-D}).
\end{itemize}

\begin{itemize}
 \item $RI$ introduces new explicit representations (e.g., your query or added assumptions).
 \item $EC$ applies algorithmic transformations (solving equations, rewriting formulas, running code).
 \item $ED$ checks or analyzes explicit content (dimension checks, algebraic consistency, syntactic
 validity).
\end{itemize}

No tacit judgment is involved here.

\subsubsection*{DHO Example}

For the DHO:

\begin{itemize}
 \item $RI$: you give the system the ODE:
 \[
 \ddot{x}(t) + 2\gamma \dot{x}(t) + \omega_0^2 x(t) = 0
 \]
 and parameters $\gamma, \omega_0$.
 \item $EC$: the system rewrites it into first-order form or finds a symbolic solution.
 \item $ED$: the system differentiates the solution and checks that it satisfies the ODE.
\end{itemize}

\subsubsection*{Example prompts}

\paragraph{$RI$.}
\begin{quote}
\textbf{User:} 
``Consider the damped harmonic oscillator
$\ddot{x}(t) + 2\gamma \dot{x}(t) + \omega_0^2 x(t) = 0$ with $\gamma = 0.2$, $\omega_0 = 1.0$.'' 
\end{quote}

\paragraph{$EC$.}
\begin{quote}
\textbf{User:} 
``Rewrite this ODE as a first-order state-space system. Do algebra only.'' 
\end{quote}

\paragraph{$ED$.}
\begin{quote}
\textbf{User:} 
``Check explicitly that the proposed solution satisfies the ODE. Show the differentiation steps.'' 
\end{quote}

\subsubsection*{Tracking notes}

\begin{itemize}
 \item $RI$, $EC$, $ED$ span all explicit machine operations.
 \item All AI-side workflows are compositions of these plus structural operators (introduced later).
\end{itemize}

\newpage

\subsection{A.0.4 Tacit Category $T$ (Companion)}

\paragraph{Reference in Appendix A.}
See Definitions~1.7--1.9.

\subsubsection*{Plain-language view}

You define the tacit category $T$ as the part of $C_{\text{full}}$ whose morphisms land in $T$.
Its generators include:

\begin{itemize}
 \item $TE : T \to T$ (Tacit Evaluation, pronounced \emph{T-E}),
 \item $HFD : T \to T$ (Human Feedback Decision, pronounced \emph{H-F-D}).
\end{itemize}

\begin{itemize}
 \item $TE$: the human evaluating content received via $RB$.
 \item $HFD$: the human deciding what to do next (approve, reject, request clarification, change the
 goal).
\end{itemize}

These do not produce explicit content directly; they change the tacit state. 
Explicit content re-enters the system only when the human expresses something new, which is modeled
as $RI$.

\subsubsection*{DHO Example}

Given a DHO solution and plots, you might:

\begin{itemize}
 \item think ``this looks physically reasonable'' ($TE$),
 \item then decide ``I want to explore the overdamped regime next'' ($HFD$),
 \item then explicitly say: ``Please recompute for $\gamma = 2\omega_0$'' (this is $RI$).
\end{itemize}

Your evaluation and decision are in $T$; the explicit instruction is in $E$.

\subsubsection*{Example prompts}

\paragraph{Tacit evaluation request.}
\begin{quote}
\textbf{System:} 
``I have derived and plotted the solution for your parameters. 
Please review and tell me whether it aligns with your expectations physically. If not, specify what you consider incorrect or surprising.'' 
\end{quote}

\paragraph{Human feedback decision.}
\begin{quote}
\textbf{User:} 
``This seems correct. Now compute the energy decay over time and show it on a log-scale plot.'' 
\end{quote}

Here, your tacit approval led to an explicit follow-up request ($RI$).

\subsubsection*{Tracking notes}

\begin{itemize}
 \item $TE$ and $HFD$ are the basic tacit morphisms.
 \item They live exclusively in the human domain $T$.
 \item Explicit content emerges again only via $RI$.
\end{itemize}

\newpage

\subsection{A.0.5 Boundary Morphism and $RB$ Functor (Companion)}

\paragraph{Reference in Appendix A.}
See Definitions~1.10--1.11 (boundary morphism and $RB^\sharp$, etc.).

\subsubsection*{Plain-language view}

You refine the role of the boundary morphism $RB : E \to T$ and introduce a functorial viewpoint:
roughly, $RB^\sharp$ (pronounced \emph{R-B-sharp}) maps explicit pipelines to tacit states that represent
``what the human has seen'' and can now evaluate.

The key idea:

\begin{itemize}
 \item There is a functor from the explicit world of computations into the tacit world of possible
 human-evaluable contents.
 \item It respects composition: longer explicit pipelines correspond to more complex inputs to tacit
 evaluation, but the mapping is well-behaved.
\end{itemize}

\subsubsection*{DHO Example}

Suppose you have an explicit pipeline:

\begin{enumerate}
 \item $RI$: introduce the DHO equation and parameters.
 \item $EC$: derive the symbolic solution.
 \item $EC$: generate a numerical plot.
 \item $ED$: verify the solution satisfies the ODE.
\end{enumerate}

Taken together, this is an explicit pipeline $P$ in $E$.

The functor $RB^\sharp$ assigns to $P$ a tacit state representing ``the content presented to the human
for evaluation'':

\begin{itemize}
 \item solution formula,
 \item error checks,
 \item plots.
\end{itemize}

Once you see and think about all this, you can apply $TE$ and $HFD$.

\subsubsection*{Example prompts}

\begin{quote}
\textbf{Prompt:} 
``Perform all explicit steps needed to solve and verify the DHO for the given parameters. 
Then present me with the solution, verification steps, and plots, but do not comment on their correctness. 
Wait for my evaluation.'' 
\end{quote}

This describes a pipeline that ends at $RB$, followed by tacit operations you perform.

\subsubsection*{Tracking notes}

\begin{itemize}
 \item $RB$ and its functorial extension link pipelines in $E$ to tacit inputs in $T$.
 \item Composition in explicit pipelines translates into structured tacit inputs.
\end{itemize}

\newpage

\subsection{A.0.6 Phase-0 and Phase-1 Axioms (Companion)}

\paragraph{Reference in Appendix A.}
See Definition~1.12 (Phase--0 constraints) and Phase--1 axioms that follow.

\subsubsection*{Plain-language view}

You formalize constraints on the \emph{operational environment} for FRFP, such as:

\begin{itemize}
 \item IL (Interaction Limitation): bounded, turn-based interaction.
 \item AR (Alignment Restriction): the system must not autonomously act beyond explicit instructions.
 \item MD (Mode Discipline): explicit vs tacit separation enforced at the interface.
 \item NTER (No Tacit Evaluation Reconstruction): the system must not reconstruct tacit evaluation 
 from explicit traces.
\end{itemize}

Phase--1 axioms add further structure to how explicit pipelines may be formed, how they may access
tools, etc., but always under ETS and the operational constraints above.

\subsubsection*{DHO Example}

For the DHO setting, IL and AR imply:

\begin{itemize}
 \item The AI waits for your explicit instructions before doing new things.
 \item It does not autonomously decide to explore unrelated parameter regimes.
\end{itemize}

MD and NTER imply:

\begin{itemize}
 \item The UI/UX never smuggles tacit evaluation back into explicit form without you typing it.
 \item The system does not mine your timing, mouse movements, or other signals to reconstruct tacit
 acceptance.
\end{itemize}

\subsubsection*{Example prompts and behaviors}

\paragraph{Phase-0 compliant behavior.}
\begin{quote}
\textbf{User:} 
``Solve the DHO for these parameters and show the explicit solution and plots.'' 

\textbf{System:} 
Performs only the requested computations, then stops and waits. 
\end{quote}

\paragraph{Non-compliant behavior (should be disallowed).}
\begin{quote}
\textbf{System (not allowed):} 
``You didn t ask, but I decided to run an experiment with different damping and take real-world actions.'' 
\end{quote}

FRFP forbids that: the system is bound to explicit instructions and the explicit--tacit separation.

\subsubsection*{Tracking notes}

\begin{itemize}
 \item Phase--0 and Phase--1 axioms constrain how FRFP systems are allowed to behave.
 \item They enforce limited interaction, mode discipline, and no tacit reconstruction.
\end{itemize}

\newpage

\section{A.1 Explicit Algebra and TDG (Companion)}

This section comments on your \textbf{A.1 Explicit Algebra and TDG}, which introduces:

\begin{itemize}
 \item TDG (Task Decomposition Grammar) signature and free category,
 \item The grammar of explicit pipelines,
 \item Shape categories,
 \item Explicit reduction systems.
\end{itemize}

\subsection{A.1.1 TDG Signature and Free Category (Companion)}

\paragraph{Reference in Appendix A.}
See Definitions~2.1--2.2 (TDG signature, free category) and Theorem~2.3.

\subsubsection*{Plain-language view}

You define a \emph{signature} for TDG (Task Decomposition Grammar) which lists basic operations like:

\begin{itemize}
 \item primitive tasks,
 \item combinators for sequencing, branching, looping,
 \item possibly tool calls.
\end{itemize}

The ``free category'' on this signature is the category whose morphisms correspond to all well-formed
compositions of these operations, with no extra equations imposed beyond the category axioms.

This free category captures all syntactically well-formed explicit pipelines as TDG terms.

\subsubsection*{DHO Example}

For the DHO, think of primitive TDG operations like:

\begin{itemize}
 \item \textsf{IntroduceODE} (introduce the DHO equation),
 \item \textsf{SolveSymbolic},
 \item \textsf{SolveNumeric},
 \item \textsf{CheckSolution},
 \item \textsf{PlotSolution}.
\end{itemize}

A full explicit pipeline is a TDG term such as:
\[
\textsf{IntroduceODE} \;\to\; \textsf{SolveSymbolic} \;\to\; \textsf{CheckSolution} \;\to\; \textsf{PlotSolution}.
\]

The free category of TDG terms has morphisms that correspond exactly to such structured pipelines.

\subsubsection*{Prompts}

\begin{quote}
\textbf{User:} 
``Give me the TDG term corresponding to: introduce the DHO, solve symbolically, verify explicitly, then plot.'' 
\end{quote}

A FRFP-aware system could respond with a symbolic TDG term in some internal syntax and use it to
drive the computation.

\newpage

\subsection{A.1.2 TDG Grammar (Companion)}

\paragraph{Reference in Appendix A.}
See the subsection on TDG grammar.

\subsubsection*{Plain-language view}

Here you introduce the grammar rules that generate valid TDG terms. This is akin to a
context-free grammar for programs, except the programs are explicit pipelines.

Non-terminals can represent:

\begin{itemize}
 \item tasks,
 \item subtasks,
 \item pipelines.
\end{itemize}

Terminals are primitive operations: $RI$, $EC$, $ED$, and specific domain operations (e.g., calling a
solver).

\subsubsection*{DHO Example}

A simplified grammar might say:

\begin{itemize}
 \item A \textsf{Pipeline} is either a \textsf{Primitive} or a
 \textsf{Primitive} followed by another \textsf{Pipeline}.
 \item A \textsf{Primitive} could be any of:
 \textsf{IntroduceODE}, \textsf{SolveSymbolic}, \textsf{SolveNumeric}, \textsf{CheckSolution},
 \textsf{PlotSolution}.
\end{itemize}

Then a well-formed TDG term for the DHO is something like:

\[
\textsf{Pipeline} \Rightarrow \textsf{IntroduceODE}\;\textsf{SolveSymbolic}\;\textsf{CheckSolution}\;\textsf{PlotSolution}.
\]

\subsubsection*{FRFP-style example}

\begin{quote}
\textbf{User:} 
``Describe the explicit pipeline to solve and verify the DHO as a TDG term, using the grammar from Appendix~A.'' 
\end{quote}

\subsubsection*{Tracking notes}

\begin{itemize}
 \item TDG is the formal language of explicit pipelines.
 \item The free category over the TDG signature captures all composition structures.
\end{itemize}

\newpage

\subsection{A.1.3 Shapes and Shape Category (Companion)}

\paragraph{Reference in Appendix A.}
See the ``Shapes and Shape Category'' subsection.

\subsubsection*{Plain-language view}

You introduce \emph{shapes} as abstract patterns of pipelines---ignoring specific content but keeping
the structural outline (branching, looping, parallelism).

The shape category $S$ (pronounced \emph{S}) has shapes as objects and shape transformations as
morphisms.

\subsubsection*{DHO Example}

Two different DHO-related pipelines may share the same shape:

\begin{itemize}
 \item Pipeline A: introduce ODE, solve symbolically, check, plot.
 \item Pipeline B: introduce ODE, solve numerically, check, plot.
\end{itemize}

Their shapes both look like:

\[
\textsf{Introduce} \to \textsf{Solve}? \to \textsf{Check} \to \textsf{Plot}
\]

where \textsf{Solve}? can be filled with different solving methods.

\subsubsection*{FRFP prompts}

\begin{quote}
\textbf{User:} 
``Abstract away from the specific solver method and give me the shape of the pipeline you used for the DHO.'' 
\end{quote}

\subsubsection*{Tracking notes}

\begin{itemize}
 \item Shapes generalize over particular pipelines.
 \item The shape category supports reasoning at a structural level.
\end{itemize}

\newpage

\subsection{A.1.4 Explicit Reduction System (Companion)}

\paragraph{Reference in Appendix A.}
See the ``Explicit Reduction System'' subsection.

\subsubsection*{Plain-language view}

You introduce a reduction system on TDG terms (pipelines), e.g., simplifying or normalizing them:

\begin{itemize}
 \item eliminating redundant steps,
 \item normalizing composition,
 \item ensuring a canonical form.
\end{itemize}

This leads to confluence and canonicality results later.

\subsubsection*{DHO Example}

Two different-looking pipelines might be semantically equivalent, e.g.:

\begin{itemize}
 \item pipeline with unnecessary double-checking or re-solving,
 \item pipeline with a minimal set of steps.
\end{itemize}

The reduction system rewrites the larger pipeline into a simpler, canonical one that has the same
effect on explicit content.

\subsubsection*{FRFP prompts}

\begin{quote}
\textbf{User:} 
``Given this TDG term for the DHO analysis, normalize it using the explicit reduction rules from Appendix~A and show me the simplified pipeline.'' 
\end{quote}

\newpage

\section{A.2 Navigation and Shapes (Companion)}

\subsection{A.2.1 Navigation Category and Groupoid (Companion)}

\paragraph{Reference in Appendix A.}
See the ``Navigation Category and Groupoid'' subsection.

\subsubsection*{Plain-language view}

You define a category (or groupoid) $N$ (pronounced \emph{N}) that captures navigation moves between
shapes or states:

\begin{itemize}
 \item moving forward/backward in a pipeline,
 \item jumping between related shapes,
 \item inverting some moves (hence groupoid).
\end{itemize}

\subsubsection*{DHO Example}

In DHO analysis, navigation might mean:

\begin{itemize}
 \item going from current explicit pipeline to a variant: e.g., switching from symbolic to numerical
 solve.
 \item moving from one parameter regime to another (underdamped $\leftrightarrow$ overdamped).
\end{itemize}

The navigation groupoid models these reversible morphisms among shapes/pipelines.

\subsubsection*{FRFP prompts}

\begin{quote}
\textbf{User:} 
``Show me how to navigate from the underdamped solution pipeline to an overdamped one while reusing as many steps as possible.'' 
\end{quote}

\newpage

\subsection{A.2.2 Fibration $p : E \to S$ (Companion)}

\paragraph{Reference in Appendix A.}
See the fibration subsection.

\subsubsection*{Plain-language view}

You define a fibration
\[
p : E \to S
\]
(pronounced \emph{p from E to S}), assigning to each explicit pipeline its shape.

Intuitively:

\begin{itemize}
 \item The base category $S$ holds abstract shapes.
 \item The total category $E$ holds concrete explicit pipelines.
 \item The fibration $p$ forgets content and keeps structure.
\end{itemize}

\subsubsection*{DHO Example}

A specific pipeline to solve and plot DHO has a shape. 
$p$ maps that pipeline to its shape:

\begin{itemize}
 \item input $\mapsto$ $S$-object: 
 \textsf{Introduce} $\to$ \textsf{Solve} $\to$ \textsf{Check} $\to$ \textsf{Plot}.
\end{itemize}

\subsubsection*{FRFP prompts}

\begin{quote}
\textbf{User:} 
``For the DHO pipeline you used, give me its shape and explain which structural choices (branching, looping) it involves.'' 
\end{quote}

\newpage

\section{A.3 Combined System: Grothendieck Construction (Companion)}

\paragraph{Reference in Appendix A.}
See the ``Grothendieck-style semidirect product $N \ltimes E$'' section.

\subsection{A.3.1 Semidirect Product $N \ltimes E$ (Companion)}

\subsubsection*{Plain-language view}

You combine navigation $N$ and explicit pipelines $E$ into a single structure:

\[
N \ltimes E
\]
(pronounced \emph{N semidirect E}).

This captures:

\begin{itemize}
 \item how navigation transforms pipelines,
 \item how pipelines sit over shapes,
 \item how actions in $N$ act on $E$ in a coherent way.
\end{itemize}

\subsubsection*{DHO Example}

For DHO, $N \ltimes E$ can encode:

\begin{itemize}
 \item a particular pipeline (solve, verify, plot),
 \item a navigation move to a related pipeline (different solver, different parameters),
 \item the combined object representing both the current state and navigation context.
\end{itemize}

\subsection{A.3.2 Rewrite System on $N \ltimes E$ (Companion)}

You then define a rewriting system on $N \ltimes E$, so navigation plus pipelines can be normalized:

\begin{itemize}
 \item eliminating redundant moves,
 \item simplifying compound navigations,
 \item yielding canonical normal forms.
\end{itemize}

\subsubsection*{FRFP example}

\begin{quote}
\textbf{User:} 
``Given the full navigation-plus-pipeline object for my DHO exploration, normalize it using the rewrite rules in Appendix~A and show me a canonical representation of where we are in the analysis.'' 
\end{quote}

\newpage

\section{A.4 Semantics and Tacit Correctness (Companion)}

\paragraph{Reference in Appendix A.}
See the ``Semantics and Tacit Correctness'' section, including the semantic functor into a tacit
preorder-enriched category.

\subsubsection*{Plain-language view}

You build a semantics mapping from explicit pipelines into a tacit correctness structure:

\begin{itemize}
 \item There is a functor (semantic interpretation) taking pipelines (or their normal forms) into
 a tacit domain that expresses what \emph{should} be correct.
 \item This functor is used to state correctness theorems:
 if the explicit pipeline is well-formed and certain conditions hold, then the induced tacit state
 satisfies a correctness property.
\end{itemize}

\subsubsection*{DHO Example}

Semantically, a DHO pipeline that:

\begin{itemize}
 \item uses correct equations,
 \item applies valid solving rules,
 \item checks results,
\end{itemize}

should map to a tacit correctness claim like:

\begin{itemize}
 \item ``Under the given assumptions, the solution and its qualitative behavior are correct.''
\end{itemize}

FRFP does not let the AI assert this tacit correctness itself; instead, the semantics lets you prove 
a theorem:

\begin{quote}
If the pipeline satisfies conditions $X$, $Y$, $Z$, then any human using it as specified will have a
path to correct tacit evaluation.
\end{quote}

\newpage

\section{A.5 Main Theorems (Companion)}

\paragraph{Reference in Appendix A.}
See the section containing the main theorems: well-formedness, confluence, canonicality,
decidability, correctness preservation, etc.

\subsubsection*{Plain-language view}

You prove results like:

\begin{itemize}
 \item \textbf{Well-formedness:} TDG terms / pipelines are well-typed in the category structure.
 \item \textbf{Confluence:} different reduction paths lead to the same normal form.
 \item \textbf{Decidability:} certain key properties are algorithmically decidable.
 \item \textbf{Correctness preservation:} explicit rewriting preserves semantic correctness.
 \item \textbf{Semantic correctness:} under your axioms, the FRFP architecture supports correct
 tacit evaluation if humans follow the protocol.
\end{itemize}

\subsubsection*{DHO Example}

For the DHO:

\begin{itemize}
 \item multiple ways to solve the ODE (characteristic polynomial, Laplace transforms, etc.) can be
 shown to be confluent at the level of solutions.
 \item pipeline normalization ensures that different explicit workflows lead to equivalent
 representations of the solution.
\end{itemize}

\subsubsection*{FRFP prompts}

\begin{quote}
\textbf{User:} 
``Explain in simple terms what the confluence theorem for TDG means for my DHO analysis. 
Why does it mean I don't have to worry about which solving path the system took, as long as it's valid?'' 
\end{quote}

\newpage

\section{A.6 Worked Examples (Companion)}

\paragraph{Reference in Appendix A.}
See ``Worked Examples'' section, including TDG term \& shape, navigation examples, etc.

\subsubsection*{Plain-language view}

You provide concrete examples of:

\begin{itemize}
 \item specific TDG terms,
 \item their shapes,
 \item navigation between them,
 \item reduction to normal forms.
\end{itemize}

\subsubsection*{DHO-flavored parallel}

You can imagine analogous DHO examples:

\begin{itemize}
 \item A full TDG term representing ``introduce DHO $\to$ solve $\to$ verify $\to$ plot''.
 \item The shape capturing its control structure.
 \item Navigation from this pipeline to one that uses a different solver.
 \item Reduction showing that two pipelines are equivalent.
\end{itemize}

\newpage

\section{A.7 Conversation Pipeline Diagram (Companion)}

\paragraph{Reference in Appendix A.}
See the conversation pipeline diagram at the end of Appendix~A.

\subsubsection*{Plain-language view}

This diagram ties everything together:

\begin{itemize}
 \item explicit pipelines in $E$ (driven by TDG),
 \item navigation moves in $N$,
 \item the fibration over shapes $S$,
 \item boundary morphism $RB$,
 \item tacit operations $TE$, $HFD$,
 \item and the overall FRFP conversation loop.
\end{itemize}

\subsubsection*{DHO Example}

For a DHO-focused session, the conversation pipeline might look like:

\begin{enumerate}
 \item $RI$: you specify the DHO task.
 \item TDG-driven $EC$/$ED$ pipeline computes and checks solutions.
 \item $RB$: results presented to you.
 \item $TE$/$HFD$: you evaluate and decide next steps.
 \item $RI$: you issue new explicit instructions (e.g., change parameters, change model).
 \item Repeat.
\end{enumerate}

\subsubsection*{FRFP-style prompts that align with the diagram}

\begin{quote}
\textbf{User:} 
``At each step, clearly separate: (1) what you are computing explicitly, (2) what you expect me to evaluate tacitly, and (3) what new explicit options I have. 
Annotate your messages accordingly.'' 
\end{quote}

A system implementing FRFP could label its messages explicitly as:

\begin{itemize}
 \item \textsf{[EXPLICIT-RESULT]},
 \item \textsf{[REQUEST-TACIT-EVAL]},
 \item \textsf{[OPTIONS-FOR-RI]}.
\end{itemize}

\newpage

\section*{Closing Notes}

This companion:

\begin{itemize}
 \item Mirrors the structure of your Appendix~A.
 \item Avoids rewriting the formal mathematics from the appendix.
 \item Provides conceptual, example-driven commentary with a running DHO example.
 \item Suggests FRFP-style prompts that match each construction.
\end{itemize}

To turn this into a tutorial or case study, you can:

\begin{itemize}
 \item Add your own DHO-specific TDG terms and diagrams.
 \item Expand the worked examples in A.6 with concrete DHO derivations.
 \item Cross-link this companion with the main appendix via hyperref.
\end{itemize}


\part*{Poster-Style Walkthrough}
\addcontentsline{toc}{part}{Poster-Style Walkthrough}
% This file is intended to be included via \input{} or \include{}.
% It contains no preamble and no document environment.

\section*{Appendix A --- Full Poster-Style Walkthrough (Dual-Flow, Detail-Preserving)}

This is a detail-preserving reorganization of the earlier full poster style Appendix companion.\par

We keep the rich blocks, callouts, and motivations from the previous version.\par

We reorganize them into the dual-flow learning sequence from the Math Primer concept map.\par

No mathematics is changed; only the teaching order is changed.\par

Two flows:\par

Content flow: explicit objects, reductions, category E\par

Structure flow: TDG/shapes, navigation, category S/N These synchronize at the fibration and proceed to the combined system, semantics, and theorems.\par

\section*{[PHASE] PHASE I --- EXPLICIT CONTENT (WHAT THE AI CAN DO)}

\subsection*{[PHASE] BLOCK C1 --- Tacit vs Explicit Worlds}

Motivation: The AI must operate entirely in an inspectable space; tacit human reasoning cannot be accessed or simulated.\par

Tacit World (Human Only): intuition, unstated assumptions, judgment.\par

Explicit World (Human + AI): formal statements, symbolic expressions, explicit computations.\par

Note: Callout: FRFP begins by drawing a hard boundary: AI stays in the explicit world.\par

\subsection*{[PHASE] BLOCK C2 --- RI, EC, ED: Atomic Explicit Objects}

Motivation: We need a controlled vocabulary for every explicit move an AI makes.\par

RI: user/human-supplied explicit representations (definitions, assumptions, givens)\par

EC: explicit computations (symbolic, algebraic, numeric transforms)\par

ED: explicit diagnostics (clarifications, warnings, notes that do not change math)\par

 Interpretation: These are the atoms of explicit reasoning.\par

\subsection*{[PHASE] BLOCK C3 --- Pipelines: Composite Explicit Objects}

Motivation: Reasoning is not atomic; it is organized into workflows.\par

A pipeline is an explicit workflow built from RI/EC/ED.\par

Pipelines are inspectable and shareable.\par

Pipelines are the primary objects the AI transforms.\par

Key point: Key Point: If the AI reasons, it does so via a pipeline.\par

\subsection*{[PHASE] BLOCK C4 --- Explicit Reduction ($\to\_\{e\}$)}

Motivation: The explicit world needs a consistent normalization system.\par

Explicit reduction:\par

simplifies/normalizes explicit content\par

distributes through workflow structure (later made precise with TDG)\par

is designed to terminate and be locally confluent\par

 Result: Explicit workflows always normalize (to stable normal forms).\par

\subsection*{[PHASE] BLOCK C5 --- Explicit Category E (Content Category)}

Motivation: By treating pipelines as categorical objects, we gain mathematical control over workflow behavior.\par

Objects: explicit pipelines (composite explicit objects)\par

Morphisms: explicit reductions + structural identities\par

Note: Clarification: Atomic objects (RI/EC/ED) admit only identity morphisms; the nontrivial morphisms act on composite pipelines.\par

Key point: Key Benefit: Composition models multi-step reasoning; identities model do nothing. \par

\subsection*{[PHASE] BLOCK C6 --- Why This Algebraic Structure Enables Proofs}

Motivation: Proofs require stable rules for how workflows transform.\par

Because E has lawful composition and identity and a reduction relation:\par

we can state and prove normalization and confluence results\par

we can compare workflows via normal forms\par

we can preserve meaning under explicit transformations\par

Note: Callout: This is what algebraic structure suitable for formal reasoning means in FRFP: pipelines behave like mathematical objects with rule-governed transformations.\par

\section*{[PHASE] PHASE II --- PURE STRUCTURE (FORM WITHOUT CONTENT)}

\subsection*{[PHASE] BLOCK S1 --- TDG: The Structural Grammar}

Motivation: We need a grammar that guarantees workflows are well-structured and manipulable.\par

TDG provides constructors (workflow form):\par

chain(Q , Q , ) -- sequential steps\par

branch(Q , Q , ) -- case splits / alternatives\par

parallel(Q , Q , ) -- independent subwork\par

loop(Q\_cond, Q\_body) -- iterative refinement\par

RI, EC, ED -- atomic leaves\par

 Interpretation: TDG is the syntax engine for building structured pipelines.\par

\subsection*{[PHASE] BLOCK S2 --- Shapes: The Structural Skeleton}

Motivation: Separate structure from content so we can reorganize workflows without altering the math.\par

A shape is a pipeline s structure with payload removed.\par

Examples:\par

chain\par

branch(chain, chain)\par

loop(chain, branch(chain, chain))\par

 Poster Metaphor: Shapes are folders (structure); pipelines are files (content).\par

\subsection*{[PHASE] BLOCK S3 --- Shape Category S}

Motivation: To reason about structure formally, shapes must live in a category.\par

Objects: shapes\par

Morphisms: structure-preserving maps (identities + composition)\par

 Outcome: Structure becomes a first-class mathematical object.\par

\section*{[PHASE] PHASE III --- NAVIGATION (STRUCTURE-ONLY CHANGE)}

\subsection*{[PHASE] BLOCK N1 --- Navigation Category and Groupoid Intuition}

Motivation: Humans frequently ask to refactor workflows; we need a safe, reversible structural editing language.\par

Navigation consists of invertible structural rewrites (examples):\par

reorder branches\par

flatten nested branches\par

regroup parallels\par

distribute loops\par

 Key Idea: Navigation changes shape only, not explicit content.\par

Note: Why groupoid? Every navigation step has an inverse (can be undone).\par

\subsection*{[PHASE] BLOCK N2 --- Why Navigation Is Not a Group}

Motivation: Navigation operations are invertible, so why not model them as a group?\par

A group describes symmetries of a single object. Navigation, however, describes invertible rewrites between many different shapes.\par

Shapes are distinct objects\par

Navigation morphisms go between shapes\par

Each rewrite has an inverse, but not a single global carrier\par

Note: Formal insight: Navigation forms a groupoid (a category where every morphism is invertible), not a group.\par

 Callout (one-sentence formal clarification): Navigation is modeled as a groupoid rather than a group because structural rewrites are invertible but act between many distinct shapes, not as symmetries of a single shape.\par

Note: Clarification --- Why groupoid and not group: A group describes invertible operations on one fixed object, whereas navigation must describe invertible rewrites between many different shapes; this multiplicity of objects forces a groupoid rather than a group.\par

\section*{[PHASE] PHASE IV --- SYNCHRONIZATION (WHERE CONTENT MEETS STRUCTURE)}

\subsection*{[PHASE] BLOCK F1 --- Fibers: All Pipelines Inside a Folder }

Motivation: Grouping pipelines by structure ensures restructuring operations know exactly where each pipeline belongs.\par

For a shape S:\par

fiber(S) = \{ e E | shape(e) = S \}\par

 Metaphor: The fiber over S is the folder contents---all pipelines with structure S.\par

Note: Callout: e fiber(shape(e)) means: pipeline e lives in the folder labeled by its shape.\par

\subsection*{[PHASE] BLOCK F2 --- Shape Fibration p : E $\to$ S}

Motivation: We need a canonical rule that assigns each pipeline to its structural folder and keeps them aligned.\par

The fibration maps each explicit pipeline to its shape:\par

p(e) = shape(e)\par

 Key Idea: Fibration is the alignment mechanism between structure and content.\par

\subsection*{[PHASE] BLOCK F3 --- Cartesian Lifting (Safe Transport)}

Motivation: When a folder (shape) moves, every file (pipeline) must move safely and canonically.\par

Given navigation $\sigma$ : S $\to$ S' and e with p(e)=S, cartesian lifting produces e' with p(e')=S':\par

 e --------$\to$ e'
 | |
 p p 
 S ---$\sigma$$\to$ S'\par

 Callout: Liftings ensure structural refactoring never corrupts explicit math.\par

 Core Safety Insight: Cartesian lifting plays the same safety role for structure/content separation that the tacit/explicit boundary plays for human/AI separation.\par

\subsection*{[PHASE] BLOCK F4 --- Navigation Action on Pipelines}

Motivation: We need one operation that applies structural rewrites directly to explicit workflows.\par

Define:\par

$\sigma \cdot e$ = e'\par

where e' is the cartesian lift of $\sigma$ applied to e.\par

Note: Clarification: Cartesian lifting states that such an e' exists and is unique; the navigation action is the explicit operation that takes e to that e'.\par

Key point: Outcome: Structure changes; content is preserved.\par

\section*{[PHASE] PHASE V --- THE COMBINED SYSTEM (STRUCTURE + CONTENT)}

\subsection*{[PHASE] BLOCK G1 --- Why Combine Navigation and Explicit Computation?}

Motivation: Real reasoning evolves both structure and content; we need one framework that tracks both.\par

Navigation handles form, explicit reduction handles math. The combined system tracks both coherently.\par

\subsection*{[PHASE] BLOCK G2 --- Grothendieck Construction / Semidirect Product N $\ltimes$ E}

Motivation: We need a single system where each step records (i) structure change and (ii) the corresponding content-following change.\par

In N $\ltimes$ E:\par

Objects: (S, e) where e lies over S\par

Morphisms: ($\sigma$, f) where $\sigma$ changes shape and f is the induced/compatible content morphism\par

Note: Key Idea: A morphism remembers both what structure changed and how content followed. \par

\subsection*{[PHASE] BLOCK G3 --- Composition and Combined Rewrite System}

Motivation: The system must remain predictable when steps are chained.\par

Composition:\par

shape part composes ($\sigma$ $\sigma$ )\par

content part composes after reindexing / lifting\par

Combined rewrite:\par

explicit reduction steps\par

navigation steps\par

mixed sequences\par

 Outcome: Mixed workflow evolution is well-defined.\par

\subsection*{[PHASE] BLOCK G4 --- Confluence (Order Doesn t Matter)}

Motivation: Humans and AI interleave edits in varying orders; results should not depend on interaction order.\par

Assumptions ensure:\par

mixed sequences converge\par

normal forms are stable (up to navigation isomorphism)\par

Key point: Practical Benefit: You can restructure while computing and still get consistent outcomes.\par

\section*{[PHASE] PHASE VI --- SEMANTICS \& TACIT CORRECTNESS}

\subsection*{[PHASE] BLOCK M1 --- Why Semantics Is Needed}

Motivation: Explicit pipelines must ultimately produce meaning interpretable by humans.\par

Semantics is the bridge from explicit computation to human-level meaning.\par

\subsection*{[PHASE] BLOCK M2 --- Semantic Interpretation \(\llbracket$ \rrbracket\)}

Motivation: We need a rule that maps an explicit pipeline to its outcome.\par

\(\llbracket$ \rrbracket\) : E $\to$ Outcome\par

Note: Callout: This explains what a pipeline means, not just what it does.\par

\subsection*{[PHASE] BLOCK M3 --- TE and HFD (Human Roles)}

Motivation: Some judgments cannot be mechanized; humans must evaluate and decide.\par

TE: tacit evaluation\par

HFD: human final decision\par

These occur outside the explicit world.\par

\subsection*{[PHASE] BLOCK M4 --- RB Functor (Boundary / Gateway)}

Motivation: We need a controlled handoff from explicit to tacit interpretation.\par

RB is the explicit-to-tacit boundary mechanism:\par

prepares outputs for human interpretation\par

prevents hidden computation at the boundary\par

 Metaphor: RB is the airlock between explicit and tacit worlds.\par

\subsection*{[PHASE] BLOCK M5 --- Navigation Invariance of Meaning}

Motivation: Refactoring structure must not change what the result means.\par

If e and e' differ only by navigation action:\par

$\sigma \cdot e$ $\cong$ e'\par

then:\par

\(\llbracket$e\rrbracket\) = \(\llbracket$e'\rrbracket\)\par

 Outcome: Structural editing is meaning-safe.\par

\subsection*{[PHASE] BLOCK M6 --- Tacit Correctness}

Motivation: Correctness ultimately lives in human interpretation.\par

A pipeline is tacit-correct when:\par

explicit normalization behaves well\par

semantic meaning is well-defined\par

TE and HFD validate it\par

Key point: Result: FRFP outputs are human-verifiable.\par

\section*{[PHASE] PHASE VII --- FRFP THEOREMS (WHAT WE CAN NOW PROVE)}

\subsection*{[PHASE] BLOCK T1 --- Main Theorems (Poster Summary)}

Motivation: The architecture exists to make strong guarantees.\par

Theorems establish:\par

Well-formedness (pipelines stay valid)\par

Explicit completeness (AI actions live in E)\par

Minimality / initiality (the explicit world is canonical)\par

Canonicality (stable normal forms)\par

Confluence (order independence)\par

Decidability (computability of reductions)\par

Correctness preservation (meaning preserved)\par

Semantic/tacit correctness (human interpretability)\par

 Outcome: FRFP supports safe, inspectable, refactorable human--AI collaboration.\par

\section*{[PHASE] PHASE VIII --- EXAMPLES \& META-REMARKS}

\subsection*{[PHASE] BLOCK X1 --- Worked Examples (Why They Matter)}

Motivation: Examples demonstrate that the machinery is usable, not just formal.\par

Examples typically show:\par

chain pipelines reducing cleanly\par

branching workflows reorganized by navigation\par

loop/branch interaction and safe liftings\par

mixed navigation + explicit reduction reaching the same normal form\par

semantics producing human-interpretable outcomes\par

\subsection*{[PHASE] BLOCK X2 --- Meta-Remarks and Implementation Notes}

Motivation: Implementation needs guidance, not just proofs.\par

Key implementation principles:\par

represent pipelines and shapes as typed, inspectable data\par

separate navigation (structure) from reduction (content)\par

enforce explicit-only computation\par

surface semantics only at the human boundary\par

\section*{FINAL BIG PICTURE (DUAL-FLOW)}

Develop explicit content and pure structure independently, synchronize them via fibration and cartesian lifting, combine them into one coherent system, interpret outcomes semantically for humans, and prove correctness and stability by construction.\par

End of Detail-Preserving Restructured Appendix A\par


\part*{Math Primer}
\addcontentsline{toc}{part}{Math Primer}
% This file is intended to be included via \input{} or \include{}.
% It contains no preamble and no document environment.

\section*{Math Primer for Appendix A}

 Math Primer for Understanding Appendix A\par

This primer introduces the minimum mathematical concepts needed to understand Appendix A of the FRFP paper.
Everything is presented in poster-style blocks, with motivation first, intuition second, and formal meaning last.\par

No prior category theory is assumed.\par

\section*{1. OBJECTS}

\subsection*{[PHASE] BLOCK 1 --- What Is an Object?}

Motivation: We need to talk about things before we can talk about transforming them.\par

An object is simply a thing we care about.\par

\n\n\begin{itemize}\n\item a number
\item a vector
\item a workflow
\item a pipeline\n\end{itemize}

Note: Key Idea: Objects are states, configurations, or entities --- not actions.\par

\subsection*{[PHASE] BLOCK 2 --- Objects in FRFP}

Motivation: The AI must reason about explicit artifacts at multiple levels of granularity.\par

In FRFP, there are two related but distinct senses in which we talk about objects:\par

Atomic objects: RI, EC, ED\par

Composite objects: explicit pipelines built from RI, EC, ED\par

Both are objects, but they play different roles.\par

Every pipeline is a thing the system can:\n\n\begin{itemize}\n\item inspect
\item transform
\item compare\n\end{itemize}

\section*{2. MORPHISMS (TRANSFORMATIONS)}

\subsection*{[PHASE] BLOCK 1 --- What Is a Morphism?}

Motivation: Reasoning is not static --- things change.\par

A morphism is an allowed transformation from one object to another.\par

Written as:\par

A $\to$ B\par

Meaning: > There is a permitted way to go from A to B. \par

\subsection*{[PHASE] BLOCK 2 --- Morphisms as Reasoning Steps}

Motivation: Reasoning is not static --- but not everything is transformable in the same way.\par

In FRFP, morphisms primarily act on composite objects (pipelines), not on atomic objects directly.\par

Pipelines change via explicit reductions\par

Atomic objects serve as stable building blocks\par

Note: Key Idea: Objects exist at multiple levels, but not all levels have interesting morphisms.\par

Formal clarification: Atomic objects (RI, EC, ED) admit only identity morphisms, while all nontrivial morphisms of the explicit category E act on composite objects (explicit pipelines) built from them.\par

Motivation: To model reasoning, we must model steps.\par

In FRFP:\n\n\begin{itemize}\n\item morphisms represent explicit reasoning steps
\item each step is inspectable and rule-governed\n\end{itemize}

Note: Objects = states of reasoning
Note: Morphisms = steps of reasoning\par

\section*{3. COMPOSITION}

\subsection*{[PHASE] BLOCK 1 --- Why Composition Matters}

Motivation: Reasoning happens in multiple steps.\par

If:\par

A $\to$ B
B $\to$ C\par

then we can compose them:\par

A $\to$ C\par

\subsection*{[PHASE] BLOCK 2 --- Composition Is Associative}

Motivation: Grouping steps should not affect outcomes.\par

Doing:\par

(A $\to$ B) $\to$ C\par

is the same as:\par

A $\to$ (B $\to$ C)\par

 This property makes long chains of reasoning manageable.\par

\section*{4. IDENTITY MORPHISMS (STRUCTURAL IDENTITIES)}

\subsection*{[PHASE] BLOCK 1 --- Doing Nothing Is Still a Step}

Motivation: Any formal system must allow no change.\par

For every object A, there is an identity morphism:\par

id\_A : A $\to$ A\par

Meaning: > Leave A exactly as it is. \par

\subsection*{[PHASE] BLOCK 2 --- Why Structural Identities Matter}

Motivation: Without identities, composition breaks.\par

Identities guarantee:\n\n\begin{itemize}\n\item pipelines can remain unchanged
\item composition laws hold
\item equivalence can be defined\n\end{itemize}

Note: In FRFP: Structural identities mean the same pipeline is still the same pipeline.\par

\section*{5. ALGEBRAIC STRUCTURE}

\subsection*{[PHASE] BLOCK 1 --- What Does Algebraic Mean Here?}

Motivation: We want to prove things about reasoning, not just run it.\par

 Algebraic structure means:\n\n\begin{itemize}\n\item objects
\item transformations
\item rules for combining transformations
\item predictable behavior\n\end{itemize}

It does not mean equations or numbers.\par

\subsection*{[PHASE] BLOCK 2 --- Why Algebraic Structure Enables Formal Reasoning}

Motivation: Proofs require stable rules.\par

Because the system has algebraic structure, we can prove:\n\n\begin{itemize}\n\item termination
\item normalization
\item equivalence
\item correctness preservation\n\end{itemize}

Key point: Without algebraic structure, none of these claims are meaningful.\par

\section*{6. CATEGORIES (THE MINIMAL FRAMEWORK)}

\subsection*{[PHASE] BLOCK 1 --- What Is a Category?}

Motivation: We want the smallest structure that supports formal reasoning.\par

A category consists of: 1. objects 2. morphisms between objects 3. composition of morphisms 4. identity morphisms\par

That s it.\par

\subsection*{[PHASE] BLOCK 2 --- Why Categories Appear in FRFP}

Motivation: FRFP needs structure, not decoration.\par

Categories allow us to:\n\n\begin{itemize}\n\item model workflows as objects
\item model reasoning as morphisms
\item guarantee compositional correctness\n\end{itemize}

 Category theory is used because it is the weakest tool that works.\par

\section*{7. GROUPS (FOR CONTRAST)}

\subsection*{[PHASE] BLOCK 1 --- What Is a Group?}

Motivation: To understand categories and navigation, it helps to see what they generalize.\par

A group has:\n\n\begin{itemize}\n\item a single underlying object
\item elements acting as transformations of that object
\item an identity element
\item inverses for every element\n\end{itemize}

Everything is reversible and acts on one thing.\par

Note: Example intuition: Rotations of a square, permutations of a fixed set.\par

\subsection*{[PHASE] BLOCK 2 --- Why FRFP Does Not Use Groups}

Motivation: Reasoning and restructuring do not act on a single fixed object.\par

Explicit reduction:\n\n\begin{itemize}\n\item simplifies
\item loses syntactic detail
\item is directional\n\end{itemize}

Navigation:\n\n\begin{itemize}\n\item is invertible
\item but moves between different shapes\n\end{itemize}

Note: Hence: FRFP needs something more general than groups.\par

\section*{8. GROUPOIDS (THE RIGHT NOTION FOR NAVIGATION)}

\subsection*{[PHASE] BLOCK 1 --- What Is a Groupoid?}

Motivation: We need reversibility without collapsing everything to one object.\par

A groupoid is:\n\n\begin{itemize}\n\item a category
\item with many objects
\item where every morphism is invertible\n\end{itemize}

Note: Intuition: A groupoid is a group with many objects. \par

\subsection*{[PHASE] BLOCK 2 --- Why Navigation Forms a Groupoid}

Motivation: Structural refactorings must be reversible but may change form.\par

In FRFP navigation:\n\n\begin{itemize}\n\item objects = shapes
\item morphisms = invertible structural rewrites
\item rewrites go between different shapes\n\end{itemize}

Each rewrite has an inverse, but there is no single global shape.\par

Note: Conclusion: Navigation is a groupoid, not a group.\par

\subsection*{[PHASE] BLOCK 3 --- One-Sentence Formal Clarification}

Navigation is modeled as a groupoid rather than a group because structural rewrites are invertible but act between many distinct shapes, not as symmetries of a single shape.\par

\section*{9. WHY CARTESIAN LIFTING IS A CORE HUMAN--AI SAFETY IDEA}

\subsection*{[PHASE] BLOCK 1 --- Why Lifting Is Needed at All}

Motivation: If structure and content are independent, changing structure should not change computation --- so why do anything?\par

The subtlety is this:\n\n\begin{itemize}\n\item structure and content are independent in meaning
\item but content is still located inside structure\n\end{itemize}

When structure changes, the meaning must stay the same, but the placement of content must be updated.\par

\subsection*{[PHASE] BLOCK 2 --- What Goes Wrong If Lifting Is Implicit}

Motivation: Hidden mechanisms break trust.\par

If lifting were implicit:\n\n\begin{itemize}\n\item the AI could reinterpret content during restructuring
\item structure changes could smuggle in computation
\item meaning preservation would be an assumption, not a theorem\n\end{itemize}

 Silent movement is silent reasoning.\par

\subsection*{[PHASE] BLOCK 3 --- What Explicit Lifting Guarantees}

Motivation: Humans need refactoring to be as safe as mathematics.\par

By making lifting:\n\n\begin{itemize}\n\item explicit (it is a visible operation)
\item canonical (there is only one correct way)
\item unique (no ambiguity)\n\end{itemize}

you get:\n\n\begin{itemize}\n\item semantic invariance (meaning does not change)
\item confluence (order of edits does not matter)
\item human trust (nothing hidden can happen)\n\end{itemize}

 Key Insight: Cartesian lifting plays the same safety role for structure/content separation that the tacit/explicit boundary plays for human/AI separation.\par

\subsection*{[PHASE] BLOCK 4 --- One-Sentence Takeaway}

Cartesian lifting exists because structure changes the placement of computation even when it must not change its meaning, and lifting is the unique, explicit rule that updates this interface safely.\par

\section*{SUPER SUMMARY --- HOW THIS PRIMER CONNECTS TO APPENDIX A}

Objects $\to$ explicit objects (atomic RI/EC/ED and composite pipelines)\par

Morphisms $\to$ explicit reasoning steps\par

Composition $\to$ multi-step reasoning\par

Identity $\to$ structural sameness\par

Algebraic structure $\to$ provable properties\par

Category $\to$ minimal framework that supports all of the above\par

 This primer provides exactly the tools needed to understand Appendix A --- nothing more, nothing less.\par

\section*{HOW TO READ SYMBOLS (A POSTER-STYLE GUIDE)}

This page explains how to read the symbols used throughout Appendix A and the math primer. Nothing new is introduced --- this is purely about interpretation.\par

\subsection*{[PHASE] BLOCK 1 --- Objects and Membership}

Symbol:\par

e E\par

Read as:\par

 e is an explicit pipeline (an object in the explicit category). \par

Informal meaning:\par

e is a concrete, inspectable workflow\par

it lives entirely in the explicit world\par

\subsection*{[PHASE] BLOCK 2 --- Morphisms (Arrows)}

Symbol:\par

e $\to$ e'\par

Read as:\par

 There is an allowed explicit transformation from pipeline e to pipeline e'. \par

Informal meaning:\par

one step (or many steps) of explicit reasoning\par

often a simplification or normalization\par

\subsection*{[PHASE] BLOCK 3 --- Explicit Reduction}

Symbol:\par

e $\to\_\{e\}$ e'\par

Read as:\par

 Pipeline e explicitly reduces to pipeline e'. \par

Informal meaning:\par

computation happened\par

expressions simplified\par

no structural reorganization yet\par

\subsection*{[PHASE] BLOCK 4 --- Identity Morphisms}

Symbol:\par

id\_e : e $\to$ e\par

Read as:\par

 The pipeline e stays exactly the same. \par

Informal meaning:\par

a formal way of saying do nothing \par

required so pipelines behave like mathematical objects\par

\subsection*{[PHASE] BLOCK 5 --- Composition of Morphisms}

Symbol:\par

e $\to$ e' $\to$ e \par

Read as:\par

 First transform e into e', then transform e' into e . \par

Informal meaning:\par

multi-step reasoning\par

chaining explicit steps\par

\subsection*{[PHASE] BLOCK 6 --- Categories}

Symbol:\par

E, S, N\par

Read as:\par

 A collection of objects together with allowed transformations between them. \par

Informal meaning in FRFP:\par

E = explicit pipelines + explicit reductions\par

S = shapes + structure-preserving maps\par

N = navigation steps (invertible structural rewrites)\par

\subsection*{[PHASE] BLOCK 7 --- Fibers}

Symbol:\par

E\_S = \{ e E | shape(e) = S \}\par

Read as:\par

 All explicit pipelines whose structure is S. \par

Informal meaning:\par

a folder (S)\par

containing all pipelines with that structure\par

\subsection*{[PHASE] BLOCK 8 --- Functions and Functors}

Symbol:\par

p : E $\to$ S\par

Read as:\par

 A rule that maps each explicit pipeline to its shape. \par

Informal meaning:\par

a consistent translation from one kind of object to another\par

preserves relationships between objects\par

\subsection*{[PHASE] BLOCK 9 --- Semantic Interpretation}

Symbol:\par

\(\llbracket$e\rrbracket\)\par

Read as:\par

 The human-interpretable meaning of pipeline e. \par

Informal meaning:\par

what the explicit computation means\par

evaluated by humans using tacit judgment\par

\subsection*{[PHASE] BLOCK 10 --- Navigation Action}

Symbol:\par

$\sigma \cdot e$ = e'\par

Read as:\par

 Apply the structural rewrite $\sigma$ to pipeline e, producing pipeline e'. \par

Informal meaning:\par

move the workflow to a new structure\par

content follows structure safely\par

[PHASE] SUPER SUMMARY --- HOW TO READ THE MATH\par

Symbols describe:\par

what exists (objects)\par

what is allowed (morphisms)\par

how steps combine (composition)\par

how structure and content relate (functions, functors)\par

how meaning is assigned (\(\llbracket$ \rrbracket\))\par

 Rule of thumb: Every symbol in Appendix A refers to an explicit, inspectable object or transformation --- never to hidden reasoning.\par

End of How to Read Symbols Poster Page\par

\section*{CONCEPT MAP --- HOW FRFP IS BUILT (DUAL-FLOW VIEW)}

This poster-style diagram shows two parallel conceptual flows --- content and structure --- and where they are formally synchronized. Read top to bottom.\par

CONTENT FLOW (EXPLICIT REASONING) STRUCTURE FLOW (WORKFLOW FORM)

 
 OBJECTS TDG / SHAPES 
 (What the AI reasons on) (Pure structure only) 
 
 Atomic: RI, EC, ED Grammar \& skeletons 
 Composite: Pipelines 
 
 
 
 
 MORPHISMS SHAPE CATEGORY S 
 (Explicit reductions $\to\_\{e\}$) (Shapes as objects) 
 
 Identity + computation Structure-preserving 
 
 
 
 
 COMPOSITION NAVIGATION N 
 (Chaining reductions) (Invertible rewrites) 
 
 e $\to$ e' $\to$ e Reorder / refactor form 
 
 
 
 
 ALGEBRAIC STRUCTURE 
 (Rule-governed behavior) 
 
 Normalization possible 
 
 
 
 
 CATEGORY E 
 (Explicit world) 
 
 Pipelines + $\to\_\{e\}$ 
 
 
 
 
 
 FIBRATION p : E $\to$ S 
 (Content structure) 
 
 Pipelines live over 
 shapes (fibers) 
 
 
 
 
 CARTESIAN LIFTING 
 (Safe synchronization) 
 
 Structure moves 
 content follows 
 
 
 
 
 COMBINED SYSTEM N $\ltimes$ E 
 (Structure + content) 
 
 Morphisms: ($\sigma$, f) 
 
 
 
 
 SEMANTICS 
 \(\llbracket$ \rrbracket\) , RB, TE, HFD 
 
 Explicit $\to$ Meaning 
 
 
 
 
 FRFP THEOREMS 
 (What we can prove) 
 
 \item Normalization 
 \item Confluence 
 \item Correctness 
 \item Semantic invariance 
 \par

[PHASE] HOW TO READ THIS DIAGRAM\par

The left column develops explicit content and computation\par

The right column develops pure structure and reorganization\par

They are intentionally separate until the fibration\par

The fibration is the first synchronization point\par

Everything below it depends on both flows\par

Key point: Key insight: Structure never changes meaning, and content never dictates structure --- they are coordinated only via principled mechanisms.\par

[PHASE] ONE-SENTENCE BIG PICTURE\par

FRFP develops content and structure independently, synchronizes them via fibration and lifting, and only then proves global correctness and semantic stability.\par

End of Concept Map Poster Page\par


\end{document}
